\section{Institutional Background}\label{sec:institutional}

Brazil's federal SME statute (Complementary Law 123/2006, amended 2014) defines micro- and small enterprises by annual gross revenue---R\$360{,}000 (US\$103{,}000) and R\$4{,}800{,}000 (US\$1.37 million) ceilings---and mandates exclusive SME tenders for items valued at R\$80{,}000 (US\$22{,}900) or less. Public buyer units (federal, state, and municipal administrative entities subject to procurement rules) may avoid the SME-only default only by filing a formal justification with the audit courts, a process that is costly, reviewable, and sanctionable. The 2014 amendment converted what had been an option into a binding default and imposed these \emph{opt-out costs}; state-level adherence to SME-only tendering for eligible items rose from roughly 13 percent to nearly 70 percent in the subsequent years.

Since 2005, all public procurement in the state of S\~ao Paulo passes through Bolsa Eletr\^onica de Compras (BEC), a centralized electronic reverse-auction platform in which registered suppliers submit sealed bids (phase 1) followed by an iterative electronic auction (phases 2 onward) in which each participant may progressively lower its offer until the auction closes. BEC handled approximately R\$13 billion (US\$3.7 billion) of procurement in 2019 alone and has cumulatively moved more than R\$105 billion (US\$30 billion) since implementation, spanning 860{,}000 purchase offers and 4.8 million item-level transactions. BEC catalogues goods into three nested tiers---group, class, and item. Medical, dental, and hospital equipment and supplies form Group 65, within which class 6531 identifies medications subject to price ceilings set by Brazil's drug-price regulator, CMED.

Group 65 was a singular exception to the 2014 rule. Between August 2014 and February 2018, the state government and the Tribunal de Contas do Estado (TCE-SP) operated under a joint interpretation that medical supplies were strategic goods requiring open competition; during that window, no Group-65 bid was restricted to SMEs. In November 2017, the Procuradoria-Geral do Estado---a separate control agency---issued a legal opinion reversing the exception on isonomy grounds: the constitutional requirement that comparable cases be treated alike. The opinion argued that no substantive feature of medical supplies justified procedural distinctions from other product groups. Effective March 2018, opt-out costs extended to Group 65; SME-only adherence within the group jumped to 43 percent on average, below the rate for other groups largely because class-6531 pharmaceuticals are oligopolistic and often fail the three-eligible-SMEs requirement for SME-only tendering.

The identifying logic of the paper rests on two features of this timing. First, the policy change was triggered by a legalistic reinterpretation of the isonomy principle at PGE-SP, not by anything happening in the Group-65 procurement market. Public documentation turns up no sign that the decision responded to price levels, competition patterns, or supplier complaints specific to Group 65; the timing was set by the bureaucratic process inside PGE-SP and by when that office chose to communicate the opinion to public buyers. Second, the reversal applied uniformly to all Group-65 items from March 2018 on, creating a discrete treatment that can be cleanly compared against the 76 product groups already under the SME-only rule. The change is therefore plausibly exogenous to the outcomes studied, and the comparison group has the simplest structure the staggered difference-in-differences literature admits: one treated cohort against never-treated controls, free of the forbidden implicit comparisons that motivate the \citet{goodmanbacon2021}, \citet{sun2021}, \citet{callaway2021}, and \citet{borusyak2024} estimators.
